%% Chapter 15: Treewidth
%% Status: skeleton -- learning goals and section outline fixed, prose to be written.

\chapter{Treewidth}
\label{ch:treewidth}

Trees are easy. Treewidth measures how far a graph is from being a tree, and
almost every problem that is easy on trees stays easy on graphs of bounded
treewidth. This is the single most useful structural parameter in the book.

\begin{goals}
  \goalitem{define tree decompositions and treewidth}
  \goalitem{determine the treewidth of trees, cycles and series--parallel graphs}
  \poolgoal{17}{show that for any tree decomposition of $G$ some bag contains every clique of $G$, and conclude that the treewidth of $G$ is at least $|C|-1$}
  \poolgoal{16}{show, by constructing a tree decomposition, that the treewidth of a $p$ by $q$ grid graph with $p \le q$ is at most $p$}
  \goalitem{design dynamic programming algorithms over a nice tree decomposition}
  \goalitem{state Courcelle's theorem and explain what it does and does not give you}
\end{goals}

\section{Tree decompositions}
\label{sec:tree-decompositions}

\todo{Write this section.}
% Planned content: Definition, the three conditions, and first examples.

\section{Bags as separators}
\label{sec:bags-separators}

\todo{Write this section.}
% Planned content: The separator property, and the fact that every clique lives inside a single bag.

\section{Treewidth of grids}
\label{sec:treewidth-grids}

\todo{Write this section.}
% Planned content: An explicit decomposition of the $p \times q$ grid of width $p$, and the matching lower bound.

\section{Nice tree decompositions}
\label{sec:nice-decompositions}

\todo{Write this section.}

\section{Dynamic programming over tree decompositions}
\label{sec:tw-dp}

\todo{Write this section.}
% Planned content: \lang{Independent Set}, \lang{Dominating Set}, and the general shape of the argument.

\section{Courcelle's theorem}
\label{sec:courcelle}

\todo{Write this section.}

\section{Computing treewidth}
\label{sec:computing-treewidth}

\todo{Write this section.}

\section{Exercises}
\label{sec:ex-treewidth}

\todo{Write this section.}
